






\appendix
\section{Proof of Lemma \ref{lem:1rank}}\label{ap:rk}

We slightly sharpen the analysis of \cite{arsovski2021padic} here. We first show $M_{p^k,1}$ has an explicit decomposition as a product of lower and upper triangular matrices.

\begin{lem}[Lemma 5 in \cite{arsovski2021padic}]\label{lem:diag}
Let $V_m$ for $m\in \N$ be an $m\times m$ matrix whose row and columns are labelled by elements in $\{0,\hdots,m-1\}$ such that its $i,j$th entry is $z^{ij}\in \Z[z]$. 

In this setting, there exists a lower triangular matrix $L_{m}$ over $\Z[z]$ with ones on the diagonal such that its inverse is also lower triangular with entries in $\Z[z]$ with ones on the diagonal, and an upper triangular matrix $D_{m}$ over $\Z[z]$ whose rows and columns are indexed by points in $\{0,\hdots,m-1\}$ such that the $j$th diagonal entry for $j\in \{0,\hdots,m-1\}$ equals

 \begin{align*}
 D_{m}(j,j)=      \prod_{i=0}^{j-1} (z^j-z^i)
 \end{align*}
 such that $V_{m}=L_{m}D_{m}$.
\end{lem}
This statement is precisely Lemma 5 in \cite{arsovski2021padic}. We will also need Lucas's theorem from \cite{Lucas}.

% We give an alternate proof for completeness.
% \begin{proof}
% Let $E_{i,j}(x)$ be the elementary matrix which on right multiplication with a matrix $M$ adds $x$ times the $i$th row to the $j$th row. If $i<j$ then we see that $E_{i,j}(x)$ is lower triangular and $E_{i,j}(x)^{-1}=E_{i,j}(-x)$.
% We prove the lemma by proving the following statement by induction: There exists a sequence of lower triangular elementary matrices $E_{i_1,j_1}(f_1),\hdots,E_{i_{n(m)},j_{n(m)}}(f_{n(m)})$ such that  $\prod_{k=1}^{n(m)} E_{i_k,j_k}(f_k) V_m=D_m$ where $D_m$ is upper triangular and its diagonal entries are as given in the lemma statement. It is clear that proving this statement also proves the Lemma.

% For $m=1$ the statement is trivial.
% Let it be true for some $m$. In $V_{m+1}$ we take its $0$th row and consider the matrix $D'_{m+1}=\prod_{j=1}^{m} E_{0,j}(-z^j) V_{m+1}$. This operation will lead to the first column of $D'_{m+1}$ being $(1,0,\hdots,0)$ and the sub-matrix of $D'_{m+1}$ after deleting the $0$th row and the $0$th column is going to be $(z-1)V_m$. We now apply the induction hypothesis to complete the proof.
% \end{proof}



\begin{thm}[Lucas's Theorem~\cite{Lucas}]\label{thm:luc}
Let $p$ be a prime and Given two natural numbers $a$ and $b$ with expansions $a_k p^k+\hdots+a_1p+a_0$ and  
$b_kp^k+\hdots+b_0$ in base $p$ we have,
\em{$$\binom{a}{b} \text{(mod }p\text{)}=\prod\limits_{i=0}^k\binom{a_i}{b_i} \text{(mod }p\text{)} .$$}
A particular consequence is that $\binom{a}{b} \mod p$ is non-zero if and only if every digit in base-$p$ of $b$ is at most as large as every digit in base-$p$ of $a$.
\end{thm}

\begin{replem}{lem:1rank}
The $\F_p$-rank of $M_{p^\ell,n}$ is at least 
{\em $$\text{rank}_{\F_p} M_{p^\ell,n} \ge \binom{\lceil p^\ell \ell^{-1}\rceil +n}{n}.$$}
\end{replem}

\begin{proof}
Using the previous lemma we note that $V_{p^\ell}=L_{p^\ell}D_{p^\ell}$. Under the ring map $f$ from $\Z[z]$ to $\Z[z]/\langle z^{p^\ell}-1\rangle$, $V_{p^\ell}=L_{p^\ell}D_{p^\ell}$ becomes $M_{p^\ell,1}=\overline{L}_{p^\ell}\overline{D}_{p^\ell}$ where $\overline{L}_{p^\ell}$ and $\overline{D}_{p^\ell}$ are the matrices $f(L_{p^\ell})$ and $f(D_{p^\ell})$ respectively.

We next notice that $M_{p^\ell,n}$ is $M_{p^\ell,1}$ tensored with itself $n$ times which we denote as $M_{p^\ell,n}=M_{p^\ell,1}^{\otimes n}$. Using $M_{p^\ell,1}=\overline{L}_{p^\ell}\overline{D}_{p^\ell}$ and Fact \ref{fact:multOfKronecker} we have $M_{p^\ell,n}=M_{p^\ell,1}^{\otimes n}=\overline{L}_{p^\ell}^{\otimes n}\overline{D}_{p^\ell}^{\otimes n}$. As $L_{p^\ell}$ was invertible with its inverse also having entries in $\Z[z]$ we see that $\overline{L}_{p^\ell}$ is also invertible and $\left(\overline{L}_{p^\ell}^{\otimes n}\right)^{-1}M_{p^\ell,n}= \overline{D}_{p^\ell}^{\otimes n}$. Using Lemma \ref{lem:crankMatMult} we have that,
$$\rank_{\F_p} M_{p^\ell,n}\ge \rank_{\F_p} \overline{D}_{p^\ell}^{\otimes n}.$$
As $\overline{D}_{p^\ell}$ is upper triangular so will $\overline{D}_{p^\ell}^{\otimes n}$ be. Therefore, to lower bound the rank of $\overline{D}_{p^\ell}^{\otimes n}$ we can lower bound the number of non-zero diagonal elements. 

The diagonal elements in $\overline{D}_{p^\ell}^{\otimes n}$ correspond to the product of diagonal elements chosen from $n$ copies of $\overline{D}_{p^\ell}$. Recall, the rows and columns of $\overline{D}_{p^\ell}$ are labelled by $j\in \{0,\hdots,p^\ell-1\}$ with
\begin{align*}
 D_{m}(j,j)=      \prod_{i=0}^{j-1} (z^j-z^i)
 \end{align*}
Setting $z-1=w$ we note that $\F_p[z]/\langle z^{p^\ell}-1\rangle$ is isomorphic to $\F_p[w]/\langle w^{p^\ell}\rangle$. $D_{p^\ell}(j,j)$ can now be written as
\begin{align*}
 D_{m}(j,j)=     (1+w)^{j(j-1)/2} \prod_{i=1}^{j} ((1+w)^i-1)
 \end{align*}
Using Lucas's Theorem (Theorem \ref{thm:luc}) we see that the largest power of $t$ which divides $(1+w)^l-1$ is the same as the largest power of $p$ which divides $l$. For any $j\le p^\ell-1$ , therefore the largest power of $w$ which divides $\overline{D}_{p^\ell}(j,j)$ is at most
\begin{align*}
    \sum\limits_{t=0}^{\lfloor \log_p(j)\rfloor} \left(\left\lfloor \frac{j}{p^t} \right\rfloor-\left\lfloor \frac{j}{p^{t+1}} \right\rfloor\right)p^t&=j+\sum\limits_{t=1}^{\lfloor \log_p(j)\rfloor} \left\lfloor \frac{j}{p^t} \right\rfloor p^{t-1}(p-1)\\
    &\le j(1+\lfloor \log_p(j)\rfloor(1-1/p))\\
    &\le j(\ell -(\ell-1)/p).\numberthis \label{eq:jBd}
\end{align*}
Consider the set of tuples $(j_1,\hdots, j_n)\in \N$ such that $j_1+\hdots+j_n\le \lceil p^\ell/\ell \rceil $. Using \eqref{eq:jBd} we see that the diagonal entry in $\overline{D}_{p^\ell}^{\otimes n}$ corresponding to the tuple will be divisible by at most $w^{\lceil p^\ell/\ell \rceil(\ell -(\ell-1)/p)}$. It is easy to check that $\lceil p^\ell/\ell \rceil(\ell -(\ell-1)/p)\le p^\ell -1$ which will guarantee that the $(j_1,\hdots,j_n)$'th diagonal entry of $\overline{D}_{p^\ell}^{\otimes n}$ is non-zero. 

This gives us at least $\binom{\lceil p^\ell \ell^{-1} \rceil+n}{n}$ non-zero diagonal entries proving the desired rank bound.
\end{proof}